%% kapa: exam=2 round=30 year=2019 period=2 subject=appraisal_theory subject_ko=감정평가이론 seq=1 number=1 points=40
%% source: Q-Net 감정평가사 기출문제 (cst00309) artlSeq=5205595 file=raw/1008/5205595/감정평가이론.pdf
%% conversion: kordoc 4.24.51 -> tools/gen2.py (auto); review=proofread-gpt-terra; flags=none
\begin{longquestion}{【문제 1】}
공기업 A는 소지를 신규취득하고 직접 조성비용을 투입하여 택지를 조성한 후, 선분양방식에 의해 주택공급을 진행하려고 하였다. 그러나 「주택공급에 관한 규칙」의 변경에 따라 후분양방식으로 주택을 공급하려고 한다. 다음의 물음에 답하시오. (40점) \par
\medskip\noindent 물음 1) 선분양방식으로 진행하려는 시점에서 A사가 조성한 택지의 감정평가방법을 설명하시오. (10점) \par
\medskip\noindent 물음 2) 상기 개발사업을 후분양방식으로 진행하면서 택지에 대한 감정평가를 실시한다고 할 경우, 최유효이용의 관점에서 감정평가방법을 제안하시오. (10점) \par
\medskip\noindent 물음 3) ‘예상되는 분양대금에서 개발비용을 공제하여 대상획지의 가치를 평가’하는 방법에서 분양대금의 현재가치 산정과 개발비용의 현재가치 산정 시 고려할 점을 설명하시오. (20점) \par
\end{longquestion}
